% SPEC-TH-031 — Skills as Files

% Shared preamble for all spec documents.
% Include from each spec: % Shared preamble for all spec documents.
% Include from each spec: % Shared preamble for all spec documents.
% Include from each spec: \input{../preamble.tex} (adjust depth).

\documentclass[11pt,a4paper]{article}

\usepackage[utf8]{inputenc}
\usepackage[T1]{fontenc}
\usepackage{lmodern}
\usepackage[a4paper,margin=2.3cm]{geometry}
\usepackage{parskip}
\usepackage{microtype}
\usepackage[dvipsnames]{xcolor}
\usepackage{enumitem}
\usepackage{listings}
\usepackage{tcolorbox}
\usepackage{titlesec}
\usepackage{fancyhdr}
\usepackage{array}
\usepackage{longtable}
\usepackage{booktabs}
\usepackage{amsmath}
\usepackage{amssymb}
\usepackage{hyperref}
\usepackage{cleveref}

\tcbuselibrary{skins,breakable}

\definecolor{specBlue}{RGB}{30,64,132}
\definecolor{specGray}{RGB}{90,90,90}
\definecolor{specAccent}{RGB}{198,40,40}
\definecolor{specGreen}{RGB}{30,110,60}

\hypersetup{
  colorlinks=true,
  linkcolor=specBlue,
  urlcolor=specBlue,
  citecolor=specBlue,
  pdfborder={0 0 0}
}

\titleformat{\section}{\Large\bfseries\color{specBlue}}{\thesection}{0.75em}{}
\titleformat{\subsection}{\large\bfseries\color{specBlue!80}}{\thesubsection}{0.6em}{}

\makeatletter
\newcommand{\specID}[1]{\def\@specID{#1}}
\newcommand{\specTitle}[1]{\def\@specTitle{#1}}
\newcommand{\specStatus}[1]{\def\@specStatus{#1}}
\newcommand{\specVersion}[1]{\def\@specVersion{#1}}
\newcommand{\specOwner}[1]{\def\@specOwner{#1}}
\newcommand{\specTraces}[1]{\def\@specTraces{#1}}

\specID{SPEC-XXX}
\specTitle{Untitled Spec}
\specStatus{DRAFT}
\specVersion{0.1}
\specOwner{Jeremie Nunez}
\specTraces{}

\newcommand{\makespecheader}{%
  \begin{center}
    {\LARGE\bfseries\color{specBlue}\@specTitle}\\[0.4em]
    {\small\color{specGray}\texttt{\@specID} \quad v\@specVersion \quad \textsc{\@specStatus} \quad \@specOwner}
  \end{center}
  \vspace{0.4em}
  \hrule
  \vspace{0.8em}
}

\pagestyle{fancy}
\fancyhf{}
\fancyhead[L]{\small\color{specGray}\@specID}
\fancyhead[R]{\small\color{specGray}\@specTitle}
\fancyfoot[C]{\small\color{specGray}\thepage}
\renewcommand{\headrulewidth}{0pt}
\makeatother

\newtcolorbox{invariant}{
  colback=specBlue!5, colframe=specBlue!75,
  title=Invariant, fonttitle=\bfseries, breakable,
  boxrule=0.5pt, arc=2pt
}

\newtcolorbox{scenario}[1]{
  colback=white, colframe=specBlue!50,
  title={Scenario: #1}, fonttitle=\bfseries, breakable,
  boxrule=0.5pt, arc=2pt
}

\newtcolorbox{ac}[1]{
  colback=specAccent!5, colframe=specAccent!60,
  title={Acceptance Criterion #1}, fonttitle=\bfseries, breakable,
  boxrule=0.5pt, arc=2pt
}

\newtcolorbox{nongoal}{
  colback=specGray!8, colframe=specGray!60,
  title=Non-Goal, fonttitle=\bfseries, breakable,
  boxrule=0.5pt, arc=2pt
}

\lstdefinestyle{codestyle}{
  basicstyle=\ttfamily\small,
  breaklines=true,
  showstringspaces=false,
  frame=single,
  framesep=4pt,
  rulecolor=\color{specBlue!30},
  backgroundcolor=\color{specBlue!2},
  xleftmargin=0.5em,
  xrightmargin=0.5em,
  keywordstyle=\color{specBlue}\bfseries,
  commentstyle=\color{specGray}\itshape,
  stringstyle=\color{specGreen}
}
\lstset{
  inputencoding=utf8,
  extendedchars=true,
  literate=
    {—}{{--}}1
    {–}{{-}}1
    {…}{{...}}1
    {’}{{'}}1
    {‘}{{'}}1
    {“}{{``}}1
    {”}{{''}}1
    {é}{{\'e}}1
    {è}{{\`e}}1
    {ê}{{\^e}}1
    {à}{{\`a}}1
    {â}{{\^a}}1
    {ç}{{\c{c}}}1
    {ô}{{\^o}}1
    {→}{{$\to$}}1
    {←}{{$\leftarrow$}}1
    {×}{{$\times$}}1
    {≥}{{$\geq$}}1
    {≤}{{$\leq$}}1
    {≠}{{$\neq$}}1
    {∈}{{$\in$}}1
    {⌈}{{$\lceil$}}1
    {⌉}{{$\rceil$}}1
    {∞}{{$\infty$}}1
    {α}{{$\alpha$}}1
    {β}{{$\beta$}}1
}
\lstdefinelanguage{TypeScript}{
  keywords={const,let,var,function,return,if,else,for,while,class,interface,type,extends,implements,import,export,from,as,async,await,new,this,true,false,null,undefined,void,any,unknown,never,string,number,boolean,readonly,private,public,protected,enum},
  sensitive=true,
  comment=[l]{//},
  morecomment=[s]{/*}{*/},
  morestring=[b]",
  morestring=[b]',
  morestring=[b]`
}
\lstset{style=codestyle}

\newcommand{\Given}{\textbf{\color{specBlue}Given} }
\newcommand{\When}{\textbf{\color{specBlue}When} }
\newcommand{\Then}{\textbf{\color{specBlue}Then} }
\providecommand{\And}{}
\renewcommand{\And}{\textbf{\color{specBlue}And} }
 (adjust depth).

\documentclass[11pt,a4paper]{article}

\usepackage[utf8]{inputenc}
\usepackage[T1]{fontenc}
\usepackage{lmodern}
\usepackage[a4paper,margin=2.3cm]{geometry}
\usepackage{parskip}
\usepackage{microtype}
\usepackage[dvipsnames]{xcolor}
\usepackage{enumitem}
\usepackage{listings}
\usepackage{tcolorbox}
\usepackage{titlesec}
\usepackage{fancyhdr}
\usepackage{array}
\usepackage{longtable}
\usepackage{booktabs}
\usepackage{amsmath}
\usepackage{amssymb}
\usepackage{hyperref}
\usepackage{cleveref}

\tcbuselibrary{skins,breakable}

\definecolor{specBlue}{RGB}{30,64,132}
\definecolor{specGray}{RGB}{90,90,90}
\definecolor{specAccent}{RGB}{198,40,40}
\definecolor{specGreen}{RGB}{30,110,60}

\hypersetup{
  colorlinks=true,
  linkcolor=specBlue,
  urlcolor=specBlue,
  citecolor=specBlue,
  pdfborder={0 0 0}
}

\titleformat{\section}{\Large\bfseries\color{specBlue}}{\thesection}{0.75em}{}
\titleformat{\subsection}{\large\bfseries\color{specBlue!80}}{\thesubsection}{0.6em}{}

\makeatletter
\newcommand{\specID}[1]{\def\@specID{#1}}
\newcommand{\specTitle}[1]{\def\@specTitle{#1}}
\newcommand{\specStatus}[1]{\def\@specStatus{#1}}
\newcommand{\specVersion}[1]{\def\@specVersion{#1}}
\newcommand{\specOwner}[1]{\def\@specOwner{#1}}
\newcommand{\specTraces}[1]{\def\@specTraces{#1}}

\specID{SPEC-XXX}
\specTitle{Untitled Spec}
\specStatus{DRAFT}
\specVersion{0.1}
\specOwner{Jeremie Nunez}
\specTraces{}

\newcommand{\makespecheader}{%
  \begin{center}
    {\LARGE\bfseries\color{specBlue}\@specTitle}\\[0.4em]
    {\small\color{specGray}\texttt{\@specID} \quad v\@specVersion \quad \textsc{\@specStatus} \quad \@specOwner}
  \end{center}
  \vspace{0.4em}
  \hrule
  \vspace{0.8em}
}

\pagestyle{fancy}
\fancyhf{}
\fancyhead[L]{\small\color{specGray}\@specID}
\fancyhead[R]{\small\color{specGray}\@specTitle}
\fancyfoot[C]{\small\color{specGray}\thepage}
\renewcommand{\headrulewidth}{0pt}
\makeatother

\newtcolorbox{invariant}{
  colback=specBlue!5, colframe=specBlue!75,
  title=Invariant, fonttitle=\bfseries, breakable,
  boxrule=0.5pt, arc=2pt
}

\newtcolorbox{scenario}[1]{
  colback=white, colframe=specBlue!50,
  title={Scenario: #1}, fonttitle=\bfseries, breakable,
  boxrule=0.5pt, arc=2pt
}

\newtcolorbox{ac}[1]{
  colback=specAccent!5, colframe=specAccent!60,
  title={Acceptance Criterion #1}, fonttitle=\bfseries, breakable,
  boxrule=0.5pt, arc=2pt
}

\newtcolorbox{nongoal}{
  colback=specGray!8, colframe=specGray!60,
  title=Non-Goal, fonttitle=\bfseries, breakable,
  boxrule=0.5pt, arc=2pt
}

\lstdefinestyle{codestyle}{
  basicstyle=\ttfamily\small,
  breaklines=true,
  showstringspaces=false,
  frame=single,
  framesep=4pt,
  rulecolor=\color{specBlue!30},
  backgroundcolor=\color{specBlue!2},
  xleftmargin=0.5em,
  xrightmargin=0.5em,
  keywordstyle=\color{specBlue}\bfseries,
  commentstyle=\color{specGray}\itshape,
  stringstyle=\color{specGreen}
}
\lstset{
  inputencoding=utf8,
  extendedchars=true,
  literate=
    {—}{{--}}1
    {–}{{-}}1
    {…}{{...}}1
    {’}{{'}}1
    {‘}{{'}}1
    {“}{{``}}1
    {”}{{''}}1
    {é}{{\'e}}1
    {è}{{\`e}}1
    {ê}{{\^e}}1
    {à}{{\`a}}1
    {â}{{\^a}}1
    {ç}{{\c{c}}}1
    {ô}{{\^o}}1
    {→}{{$\to$}}1
    {←}{{$\leftarrow$}}1
    {×}{{$\times$}}1
    {≥}{{$\geq$}}1
    {≤}{{$\leq$}}1
    {≠}{{$\neq$}}1
    {∈}{{$\in$}}1
    {⌈}{{$\lceil$}}1
    {⌉}{{$\rceil$}}1
    {∞}{{$\infty$}}1
    {α}{{$\alpha$}}1
    {β}{{$\beta$}}1
}
\lstdefinelanguage{TypeScript}{
  keywords={const,let,var,function,return,if,else,for,while,class,interface,type,extends,implements,import,export,from,as,async,await,new,this,true,false,null,undefined,void,any,unknown,never,string,number,boolean,readonly,private,public,protected,enum},
  sensitive=true,
  comment=[l]{//},
  morecomment=[s]{/*}{*/},
  morestring=[b]",
  morestring=[b]',
  morestring=[b]`
}
\lstset{style=codestyle}

\newcommand{\Given}{\textbf{\color{specBlue}Given} }
\newcommand{\When}{\textbf{\color{specBlue}When} }
\newcommand{\Then}{\textbf{\color{specBlue}Then} }
\providecommand{\And}{}
\renewcommand{\And}{\textbf{\color{specBlue}And} }
 (adjust depth).

\documentclass[11pt,a4paper]{article}

\usepackage[utf8]{inputenc}
\usepackage[T1]{fontenc}
\usepackage{lmodern}
\usepackage[a4paper,margin=2.3cm]{geometry}
\usepackage{parskip}
\usepackage{microtype}
\usepackage[dvipsnames]{xcolor}
\usepackage{enumitem}
\usepackage{listings}
\usepackage{tcolorbox}
\usepackage{titlesec}
\usepackage{fancyhdr}
\usepackage{array}
\usepackage{longtable}
\usepackage{booktabs}
\usepackage{amsmath}
\usepackage{amssymb}
\usepackage{hyperref}
\usepackage{cleveref}

\tcbuselibrary{skins,breakable}

\definecolor{specBlue}{RGB}{30,64,132}
\definecolor{specGray}{RGB}{90,90,90}
\definecolor{specAccent}{RGB}{198,40,40}
\definecolor{specGreen}{RGB}{30,110,60}

\hypersetup{
  colorlinks=true,
  linkcolor=specBlue,
  urlcolor=specBlue,
  citecolor=specBlue,
  pdfborder={0 0 0}
}

\titleformat{\section}{\Large\bfseries\color{specBlue}}{\thesection}{0.75em}{}
\titleformat{\subsection}{\large\bfseries\color{specBlue!80}}{\thesubsection}{0.6em}{}

\makeatletter
\newcommand{\specID}[1]{\def\@specID{#1}}
\newcommand{\specTitle}[1]{\def\@specTitle{#1}}
\newcommand{\specStatus}[1]{\def\@specStatus{#1}}
\newcommand{\specVersion}[1]{\def\@specVersion{#1}}
\newcommand{\specOwner}[1]{\def\@specOwner{#1}}
\newcommand{\specTraces}[1]{\def\@specTraces{#1}}

\specID{SPEC-XXX}
\specTitle{Untitled Spec}
\specStatus{DRAFT}
\specVersion{0.1}
\specOwner{Jeremie Nunez}
\specTraces{}

\newcommand{\makespecheader}{%
  \begin{center}
    {\LARGE\bfseries\color{specBlue}\@specTitle}\\[0.4em]
    {\small\color{specGray}\texttt{\@specID} \quad v\@specVersion \quad \textsc{\@specStatus} \quad \@specOwner}
  \end{center}
  \vspace{0.4em}
  \hrule
  \vspace{0.8em}
}

\pagestyle{fancy}
\fancyhf{}
\fancyhead[L]{\small\color{specGray}\@specID}
\fancyhead[R]{\small\color{specGray}\@specTitle}
\fancyfoot[C]{\small\color{specGray}\thepage}
\renewcommand{\headrulewidth}{0pt}
\makeatother

\newtcolorbox{invariant}{
  colback=specBlue!5, colframe=specBlue!75,
  title=Invariant, fonttitle=\bfseries, breakable,
  boxrule=0.5pt, arc=2pt
}

\newtcolorbox{scenario}[1]{
  colback=white, colframe=specBlue!50,
  title={Scenario: #1}, fonttitle=\bfseries, breakable,
  boxrule=0.5pt, arc=2pt
}

\newtcolorbox{ac}[1]{
  colback=specAccent!5, colframe=specAccent!60,
  title={Acceptance Criterion #1}, fonttitle=\bfseries, breakable,
  boxrule=0.5pt, arc=2pt
}

\newtcolorbox{nongoal}{
  colback=specGray!8, colframe=specGray!60,
  title=Non-Goal, fonttitle=\bfseries, breakable,
  boxrule=0.5pt, arc=2pt
}

\lstdefinestyle{codestyle}{
  basicstyle=\ttfamily\small,
  breaklines=true,
  showstringspaces=false,
  frame=single,
  framesep=4pt,
  rulecolor=\color{specBlue!30},
  backgroundcolor=\color{specBlue!2},
  xleftmargin=0.5em,
  xrightmargin=0.5em,
  keywordstyle=\color{specBlue}\bfseries,
  commentstyle=\color{specGray}\itshape,
  stringstyle=\color{specGreen}
}
\lstset{
  inputencoding=utf8,
  extendedchars=true,
  literate=
    {—}{{--}}1
    {–}{{-}}1
    {…}{{...}}1
    {’}{{'}}1
    {‘}{{'}}1
    {“}{{``}}1
    {”}{{''}}1
    {é}{{\'e}}1
    {è}{{\`e}}1
    {ê}{{\^e}}1
    {à}{{\`a}}1
    {â}{{\^a}}1
    {ç}{{\c{c}}}1
    {ô}{{\^o}}1
    {→}{{$\to$}}1
    {←}{{$\leftarrow$}}1
    {×}{{$\times$}}1
    {≥}{{$\geq$}}1
    {≤}{{$\leq$}}1
    {≠}{{$\neq$}}1
    {∈}{{$\in$}}1
    {⌈}{{$\lceil$}}1
    {⌉}{{$\rceil$}}1
    {∞}{{$\infty$}}1
    {α}{{$\alpha$}}1
    {β}{{$\beta$}}1
}
\lstdefinelanguage{TypeScript}{
  keywords={const,let,var,function,return,if,else,for,while,class,interface,type,extends,implements,import,export,from,as,async,await,new,this,true,false,null,undefined,void,any,unknown,never,string,number,boolean,readonly,private,public,protected,enum},
  sensitive=true,
  comment=[l]{//},
  morecomment=[s]{/*}{*/},
  morestring=[b]",
  morestring=[b]',
  morestring=[b]`
}
\lstset{style=codestyle}

\newcommand{\Given}{\textbf{\color{specBlue}Given} }
\newcommand{\When}{\textbf{\color{specBlue}When} }
\newcommand{\Then}{\textbf{\color{specBlue}Then} }
\providecommand{\And}{}
\renewcommand{\And}{\textbf{\color{specBlue}And} }


\specID{SPEC-TH-031}
\specTitle{Skills as Files --- versioned markdown prompts as code artifacts}
\specStatus{DRAFT}
\specVersion{0.1}
\specOwner{Jeremie Nunez}

\begin{document}
\makespecheader

\section{Context}
Thalamus cortices are driven by LLM system prompts. Embedding those prompts in TypeScript
string literals makes them invisible to non-engineers, impossible to diff meaningfully, and
awkward to review. This spec promotes prompts to first-class, versioned artifacts:
markdown files under \texttt{packages/thalamus/src/cortices/skills/}, loaded at runtime by
\texttt{cortex-llm.ts}, validated in CI, and reviewable by analysts and compliance without
touching TypeScript.

\section{Goals}
\begin{itemize}[leftmargin=1.2em]
  \item Every cortex prompt lives in a dedicated \texttt{.md} file next to code.
  \item Each skill file carries structured front-matter: \texttt{name}, \texttt{description},
        \texttt{sqlHelper}, \texttt{params}.
  \item The body declares \emph{purpose}, \emph{inputs}, \emph{outputs}, and
        \emph{examples} in fixed sections.
  \item CI validates the format so a malformed skill cannot reach \texttt{main}.
  \item Non-engineers can open a PR editing only markdown and ship a behavior change safely.
\end{itemize}

\section{Non-Goals}
\begin{nongoal}
This spec does not define the DAG execution order, the transport to the LLM provider, or
the structure of the knowledge graph writes (covered in SPEC-TH-030). It does not replace
secret configuration --- API keys remain in the environment.
\end{nongoal}

\section{Invariants}
\begin{invariant}
A skill file is the single source of truth for its cortex's prompt. The TypeScript code
MUST load the file at runtime (or at build time via a typed registry) and MUST NOT duplicate
the prompt as a string literal.
\end{invariant}

\begin{invariant}
Every skill file satisfies the schema: YAML front-matter with
\texttt{name, description, sqlHelper, params}, followed by a body that contains at least
the sections \texttt{Purpose}, \texttt{Inputs}, \texttt{Output Format}, and \texttt{Examples}.
Any deviation fails CI.
\end{invariant}

\begin{invariant}
The \texttt{Output Format} section declares a JSON shape. The runtime parses LLM responses
against that shape; rejections are logged with the skill name and a sample. No silent
coercion.
\end{invariant}

\section{Interface}

\textbf{File layout} --- a skill is a single markdown document:

\begin{lstlisting}[language=TypeScript]
---
name: strategist              # required, kebab-case, matches filename
description: >                # required, one-paragraph summary
  Meta-synthesis cortex that reads findings from upstream cortices
  and produces portfolio-level recommendations.
sqlHelper: none               # required, a registered helper or "none"
params:                       # required, may be empty
  maxFindings: number
  lang: "fr" | "en"
---

# Strategist

## Purpose
One paragraph. What does this cortex decide, and why does it exist.

## Inputs
Bulleted list of what arrives in the DATA payload.

## Output Format
Return JSON: `{ "findings": [...] }`. Each finding has fields
title, summary, findingType, confidence, impactScore, evidence, edges.

## Examples
At least one worked example: input shape -> expected finding.
\end{lstlisting}

\textbf{Loader API} --- consumed by \texttt{cortex-llm.ts}:

\begin{lstlisting}[language=TypeScript]
// packages/thalamus/src/cortices/skills/registry.ts

export interface SkillFile {
  name: string;
  description: string;
  sqlHelper: string;          // key into sql-helpers registry, or "none"
  params: Record<string, string>;
  systemPrompt: string;       // the body, minus front-matter
  sourcePath: string;         // for logs and traceability
  sha256: string;             // content hash, stamped into provenance
}

export function loadSkill(name: string): SkillFile;
export function listSkills(): SkillFile[];
export function validateSkill(raw: string): {
  ok: boolean; errors: string[];
};
\end{lstlisting}

\section{Scenarios}

\begin{scenario}{Analyst edits a skill}
\Given an analyst updates the \texttt{Output Format} section of \texttt{strategist.md} \\
\When the PR is opened \\
\Then CI parses the front-matter and validates the required sections \\
\And the diff is human-readable (markdown, not escaped strings) \\
\And no TypeScript file needs to change for the update to ship.
\end{scenario}

\begin{scenario}{Missing required section fails CI}
\Given a skill file lacks an \texttt{Output Format} section \\
\When \texttt{validateSkill} runs in CI \\
\Then the build fails \\
\And the error message names the file and the missing section.
\end{scenario}

\begin{scenario}{Filename and \texttt{name} must match}
\Given a file \texttt{strategist.md} declares \texttt{name: strategic-synth} in front-matter \\
\When the validator runs \\
\Then the build fails with \texttt{skill name must equal filename stem}.
\end{scenario}

\begin{scenario}{Runtime loads the file, not a literal}
\Given \texttt{cortex-llm.ts} needs the strategist prompt \\
\When \texttt{analyzeCortexData} is invoked \\
\Then the system prompt is loaded via \texttt{loadSkill("strategist")} \\
\And the finding's provenance records the skill's \texttt{sha256}.
\end{scenario}

\begin{scenario}{Prompt hash reaches the knowledge graph}
\Given a cortex writes a finding to the graph (see SPEC-TH-030) \\
\When the write lands \\
\Then the row's provenance includes the \texttt{skillSha256} that produced it \\
\And a later audit can diff the exact prompt version used at write time.
\end{scenario}

\section{Acceptance Criteria}

\begin{ac}{AC-1}
\texttt{validateSkill} applied to every file under
\texttt{packages/thalamus/src/cortices/skills/} returns \texttt{ok: true} on a clean tree.
\end{ac}

\begin{ac}{AC-2}
Seeding a malformed skill (missing \texttt{Output Format}) causes CI to fail with a message
naming the file and the missing section.
\end{ac}

\begin{ac}{AC-3}
\texttt{listSkills()} returns exactly the set of files in the skills directory; renaming
a file and not updating \texttt{name} in front-matter fails validation.
\end{ac}

\begin{ac}{AC-4}
No TypeScript file under \texttt{packages/thalamus/src/cortices/} contains a multi-line
string literal longer than 200 characters that looks like a system prompt; a lint rule
enforces this.
\end{ac}

\begin{ac}{AC-5}
The \texttt{sha256} of each loaded skill is deterministic across machines (normalized line
endings) and is included in provenance records.
\end{ac}

\begin{ac}{AC-6}
An analyst-only PR (markdown-only diff) passes CI and deploys without requiring engineering
sign-off beyond normal review.
\end{ac}

\begin{ac}{AC-7}
When the LLM response fails to match the declared \texttt{Output Format}, the error log
includes \texttt{skillName}, \texttt{skillSha256}, and a truncated sample of the response.
\end{ac}

\section{Traceability}

\begin{center}
\begin{tabular}{@{}lll@{}}
\toprule
AC & Test file & Test name \\
\midrule
AC-1 & \texttt{packages/thalamus/tests/skills/validate.spec.ts} & \texttt{all shipped skills are valid} \\
AC-2 & \texttt{packages/thalamus/tests/skills/validate.spec.ts} & \texttt{missing Output Format fails} \\
AC-3 & \texttt{packages/thalamus/tests/skills/registry.spec.ts} & \texttt{filename must match name} \\
AC-4 & \texttt{packages/thalamus/tests/lint/no-inline-prompts.spec.ts} & \texttt{no long prompt literals} \\
AC-5 & \texttt{packages/thalamus/tests/skills/registry.spec.ts} & \texttt{sha256 is deterministic} \\
AC-6 & \texttt{packages/thalamus/tests/ci/analyst-pr.spec.ts} & \texttt{markdown-only PR passes} \\
AC-7 & \texttt{packages/thalamus/tests/cortex-llm/output-format.spec.ts} & \texttt{mismatch is logged with sha} \\
\bottomrule
\end{tabular}
\end{center}

\section{Open Questions}
\begin{itemize}[leftmargin=1.2em]
  \item Should \texttt{params} use a JSON Schema or a lightweight DSL? JSON Schema is
        standard but verbose in front-matter.
  \item Do we version skills by semver in the front-matter as well as by git history?
  \item Should the \texttt{Output Format} block be executable (e.g. an embedded Zod schema)
        or remain prose plus a documented JSON shape?
  \item Localization --- one file per language, or a single file with tagged sections?
\end{itemize}

\end{document}
